\section{Flowed coalescence projectors}
\label{app:flowed}

Tier 4 asks a question Tiers 1--3 deliberately avoid: not \emph{how
many} states with given quantum numbers exist in the out-state, but
whether the quarks inside them are \emph{compactly} arranged.  Any
such question requires an operator that resolves sub-hadronic
structure, and bare multi-local quark operators
$\bar\psi\Gamma\psi\,\bar\psi\Gamma'\psi(x)$ mix under
renormalization with the full tower of same-quantum-number operators;
their matrix elements have no scheme-independent continuum limit.

The gradient flow provides the controlled compromise.  Fields at flow
time $t_{\rm fl}$ are smeared over the physical radius
$r_{\rm fl}=\sqrt{8t_{\rm fl}}$, and correlation functions of flowed
fields are finite as $a\to0$ at fixed $t_{\rm fl}$: for the gauge
sector this is L\"uscher's theorem \cite{Luscher:2010iy}, and flowed
fermion bilinears require only the single multiplicative wave-function
factor of Ref.~\cite{Makino:2014taa}, fixed nonperturbatively by a
vacuum normalization condition.  Define, on a certified cluster region
$R$, the flowed compactness projector
\begin{equation}
\Pi^{(4q)}_{t_{\rm fl}}(R) \;=\;
\int_R \prod_{k=1}^{4} d^3x_k \;
\Theta\!\big(\{x\};r_{\rm fl}\big)\,
\mathcal O^\dagger_{4q,t_{\rm fl}}(\{x\})\,
\mathcal O_{4q,t_{\rm fl}}(\{x\}),
\end{equation}
with $\mathcal O_{4q,t_{\rm fl}}$ a gauge-invariant four-quark
interpolator built from flowed fields and $\Theta$ a window
restricting the mutual separations to $\lesssim r_{\rm fl}$.  Its
expectation value on the cluster is finite in the continuum limit and
depends on exactly one scheme label, the physical flow radius --- a
\emph{labeled} scheme, in contrast to the unlabeled divergences of the
bare operator.

Three uses follow.  (i)~At fixed $t_{\rm fl}$, \emph{ratios} of
$\langle\Pi^{(4q)}\rangle$ across clusters of the same quantum numbers
compare compactness between states with common normalization.
(ii)~The $t_{\rm fl}$-dependence itself is the diagnostic: a compact
tetraquark saturates $\langle\Pi\rangle$ once $r_{\rm fl}$ exceeds its
core size, while a molecular state grows until $r_{\rm fl}$ reaches
the hadron-pair separation --- the operator-level version of the
coalescence yield systematics that the ExHIC program exploits
phenomenologically \cite{ExHIC:2010gcb,ExHIC:2017smd}.  (iii)~The map
$r_0 \leftrightarrow r_{\rm fl}$ connects this tier quantitatively to
the classical binding criterion of Sec.~\ref{sec:cluster-id} and,
through it, to Lebed's nearest-neighbor rule \cite{Lebed:2016jvq}:
the classical model is, in this precise sense, the
$\hbar\to0$ caricature of the Tier-4 measurement at fixed scheme
label.

What is \emph{not} claimed: no infrared-safe, scheme-independent
``compactness number'' exists, because compactness is not an
S-matrix property.  Tier 4 outputs are honest functions of
$t_{\rm fl}$, reported as such, and their value lies in
\emph{differential} statements --- between states, between system
sizes, between preparations --- where the scheme label cancels in the
comparison.